\section{Current Implementation Measurements and Discussion}
\label{sec:results}
\begin{table}[!t]
\centering
\caption{Intent Classifier: Measured Per-Class Metrics on the Reconstructed
$34$-Example Validation Split (Seed $42$). Overall Accuracy $0.5294$;
Macro $F_1$ $0.4042$; Weighted $F_1$ $0.4869$.}
\label{tab:lstmperclass}
\footnotesize
\renewcommand{\arraystretch}{1.12}
\begin{tabular}{@{}lrrrr@{}}
\toprule
\textbf{Class} & \textbf{Prec.} & \textbf{Rec.} & \textbf{$F_1$} & \textbf{Supp.} \\
\midrule
ADMISSIONS          & 1.000 & 0.200 & 0.333 & 5 \\
COURSES             & 0.400 & 0.500 & 0.444 & 4 \\
DEPARTMENTS         & 0.667 & 1.000 & 0.800 & 2 \\
ACADEMICS           & 0.667 & 1.000 & 0.800 & 2 \\
FACILITIES          & 0.600 & 1.000 & 0.750 & 3 \\
ROOM\_LOCATION      & 1.000 & 0.857 & 0.923 & 7 \\
DEPARTMENT\_LOCATION & 0.250 & 1.000 & 0.400 & 1 \\
LABORATORY\_LOCATION & 0.000 & 0.000 & 0.000 & 3 \\
VIRTUAL\_TOUR       & 0.333 & 0.500 & 0.400 & 2 \\
CAMPUS\_INFO        & 0.000 & 0.000 & 0.000 & 1 \\
GENERAL            & 0.000 & 0.000 & 0.000 & 1 \\
UNKNOWN            & 0.000 & 0.000 & 0.000 & 3 \\
\bottomrule
\end{tabular}
\end{table}

\begin{figure}[!t]
\centering
\footnotesize
\setlength{\tabcolsep}{2.2pt}
\renewcommand{\arraystretch}{1.05}
\begin{tabular}{@{}l|cccccccccccc@{}}
\multicolumn{1}{c}{} & \multicolumn{12}{c}{\scriptsize predicted} \\
\scriptsize true & \rotatebox{90}{ADM} & \rotatebox{90}{COU} & \rotatebox{90}{DEP} & \rotatebox{90}{ACA} & \rotatebox{90}{FAC} & \rotatebox{90}{ROOM} & \rotatebox{90}{DLOC} & \rotatebox{90}{LLOC} & \rotatebox{90}{VT} & \rotatebox{90}{CINF} & \rotatebox{90}{GEN} & \rotatebox{90}{UNK} \\
\hline
ADM   & \cellcolor{black!15}1 & 1 & . & . & 1 & . & . & . & 1 & 1 & . & . \\
COU   & . & \cellcolor{black!15}2 & 1 & 1 & . & . & . & . & . & . & . & . \\
DEP   & . & . & \cellcolor{black!15}2 & . & . & . & . & . & . & . & . & . \\
ACA   & . & . & . & \cellcolor{black!15}2 & . & . & . & . & . & . & . & . \\
FAC   & . & . & . & . & \cellcolor{black!15}3 & . & . & . & . & . & . & . \\
ROOM  & . & . & . & . & 1 & \cellcolor{black!15}6 & . & . & . & . & . & . \\
DLOC  & . & . & . & . & . & . & \cellcolor{black!15}1 & . & . & . & . & . \\
LLOC  & . & . & . & . & . & . & 3 & \cellcolor{black!15}0 & . & . & . & . \\
VT    & . & . & . & . & . & . & . & . & \cellcolor{black!15}1 & . & . & 1 \\
CINF  & . & 1 & . & . & . & . & . & . & . & \cellcolor{black!15}0 & . & . \\
GEN   & . & 1 & . & . & . & . & . & . & . & . & \cellcolor{black!15}0 & . \\
UNK   & . & . & . & . & . & . & . & . & 1 & 2 & . & \cellcolor{black!15}0 \\
\end{tabular}
\caption{Intent classifier confusion matrix (raw counts) on the $34$-example
validation split. Diagonal shaded. The classifier collapses
\code{LABORATORY\_LOCATION} onto \code{DEPARTMENT\_LOCATION} and scatters
\code{ADMISSIONS} and the low-support classes, consistent with overfitting
on $168$ synthetic training utterances.}
\label{fig:confusion}
\end{figure}


This section reports only measurements that actually exist in the artifact.
They are of two kinds: implementation-scale figures
(Table~\ref{tab:counts}), and the few quantitative observations captured
during development. None constitutes an experimental performance evaluation
of the kind specified in Section~\ref{sec:eval}.

\subsection{Implementation Scale}
The system integrates a $1{,}507$-chunk knowledge base (from $128$ PDFs and
$42$ web pages; $186$ URLs attempted with typed failure logging), a
$180$-node/$325$-edge campus graph, $161$ panorama scenes with $849$
cross-floor sightline hotspots across a six-bucket floor sequence, a
$12$-class LSTM over a $283$-word vocabulary, and a relational schema with
$14$ migrations. The chat tables have accumulated $378$ sessions and
$1{,}400$ messages during development, evidence that the full front-end
$\to$ API $\to$ supervisor $\to$ retrieval $\to$ Ollama $\to$ persistence
path runs end to end. \emph{These are scale figures, not accuracy figures};
in particular the chunk count and the message count say nothing about
answer quality.

\subsection{Intent Classifier}
The only quantitative machine-learning result in the artifact is the intent
classifier's held-out accuracy. Training reports \textbf{training accuracy
$1.000$} and \textbf{validation accuracy $0.529$} on a seeded $80/20$ split
of $168$ synthetic utterances ($134$ train, $34$ validation). An
independent re-evaluation, reconstructing the identical split, reproduced
accuracy $0.5294$ exactly and additionally reports macro precision $0.410$,
macro recall $0.505$, macro $F_1$ $0.404$, and weighted $F_1$ $0.487$.
Table~\ref{tab:lstmperclass} gives the per-class breakdown and
Fig.~\ref{fig:confusion} the confusion matrix.

\emph{Discussion.} The gap between training accuracy $1.00$ and validation
accuracy $0.53$ on $168$ examples is a clear overfitting
signature~\cite{goodfellow2016deep}; the model's own metadata and its
dataset module say so explicitly. Four of twelve classes
(\code{LABORATORY\_LOCATION}, \code{CAMPUS\_INFO}, \code{GENERAL},
\code{UNKNOWN}) score $F_1 = 0$ on the validation split, several with
support $1$--$3$. \code{LABORATORY\_LOCATION} is entirely absorbed into
\code{DEPARTMENT\_LOCATION}. \textbf{This classifier must not be presented
as accurate}, and the architecture does not depend on it being so: it is a
fallback signal, reached only when the deterministic router finds no match,
and accepted only above a $0.6$ probability. The high-support, lexically
distinctive classes it does need to handle in that fallback role ---
\code{ROOM\_LOCATION} ($F_1 = 0.923$), \code{DEPARTMENTS}/\code{ACADEMICS}/
\code{FACILITIES} ($F_1 \ge 0.75$) --- are its strongest. Experiment E3
(Section~\ref{sec:eval}) is what would establish real generalization, on
real logged queries.

\subsection{Confidence Thresholds}
The category thresholds $\theta_{\text{HIGH}} = 0.58$ and
$\theta_{\text{MED}} = 0.48$ were placed in the gaps observed on a
\textbf{nine-query} development sample (Table~\ref{tab:confsample}): all
three out-of-domain queries scored below $0.48$ and all six in-domain
queries scored $0.48$ or above, with a secondary gap between $0.516$ and
$0.641$ separating ``one strong supporting chunk'' from ``relevant but
thin, spread-out evidence.'' \emph{This is a calibration artifact with
$n = 9$, not a validated decision boundary.} Its behavior on a
statistically sized set is unknown and is the subject of Experiment E5. We
also note again that the intent term in Eq.~\eqref{eq:conf} is currently
the retrieval fallback, so the deployed confidence score is effectively a
function of retrieval quality alone.

\subsection{Concurrency Observation}
During a Phase-10 timeout investigation, two chat requests issued
simultaneously against the CPU-only development host were measured
completing in \textbf{$17$\,s and $45$\,s} respectively, versus an
estimated $\approx 8$\,s each in isolation --- a super-linear degradation
attributed to CPU contention between two concurrent Ollama generations
running in Starlette's thread pool with no concurrency bound. The response
was to serialize Ollama calls behind a semaphore (default width $1$), which
converts ``$N$ requests all degrade together'' into ``requests queue and
each receives the full $\approx 8$\,s service time.'' \textbf{This is a
single anecdotal engineering observation, not a load-test result}; no
throughput curve, latency percentile, or concurrency sweep exists.
Experiment E6 and a proper load test (asyncio + httpx, $\ge 50$ concurrent
sessions) are required to characterize behavior under load.

\subsection{What Is Not Measured}
For completeness and to prevent misreading: there is no retrieval
precision/recall, no MRR or nDCG, no reranker A/B comparison, no
faithfulness or answer-relevancy score, no citation-correctness rate, no
hallucination rate, no confidence-threshold sensitivity analysis, no
A$^{\star}$ latency or expanded-node benchmark, no cross-floor transition
success rate, and no usability or user-study data. Section~\ref{sec:eval}
specifies how each would be obtained.
